\section{The Skill Economy: Open Problems}
\label{sec:open}

The arrival of a commercial marketplace for robot skills---Unitree's UniStore ships
one-tap, cross-model motion packages \citep{unistore2026}---turns several research
questions from hypothetical into pressing. We frame them as the open agenda of this survey.

\paragraph{The static-to-adaptive gap.}
Every skill shipped today is, in effect, replayed: a motion package or a fixed program is
installed and executed without perceiving whether \emph{this} kitchen, gripper, or object
differs from the one it was authored for. The techniques of \S\ref{sec:code}---within-task
repair, persistent memory, and search---are precisely what would let a downloaded skill
\emph{adapt} on the target robot. Closing this gap is the difference between a store of
animations and a store of capabilities, and it is the direct application of the §3.1.5
loop to a distribution setting where the base skill is given rather than discovered.

\paragraph{Cross-embodiment portability.}
UniStore advertises a single skill running across the G1 and H1 humanoids and the B2/Go2
quadrupeds---bodies with different morphologies, action spaces, and dynamics. This is the
cross-embodiment transfer problem (\S\ref{sec:transfer}) at deployment scale: without a
shared action interface \citep{openx} or explicit embodiment adaptation, a package tuned
for one platform has no guarantee of correctness on another. Formal portability---what must
hold for a skill to be certified on a new body---remains open.

\paragraph{Provenance and trust.}
A marketplace invites third-party, crowdsourced uploads, including raw motion-capture data
and code from unknown authors. Robot skills act on the physical world, so provenance
(who authored a skill, from what data, validated how) becomes a safety property, not merely
a metadata nicety. There is no accepted standard for signing, attesting, or auditing a
robot skill's origin.

\paragraph{Safety verification.}
Vendors promise that ``every skill is scanned, tested, and verified,'' but for a
code-as-policy skill this is undecidable in general and hard in practice: the skill's effect
depends on the scene it meets. Practical questions include which pre-conditions a skill must
declare, what sandboxed or simulated screening it must pass, and how run-time monitors
(cf.\ \emph{Code-as-Monitor} \citep{codeasmonitor}) can bound behaviour on unseen inputs.

\paragraph{Skill composition.}
Downloading two skills should ideally yield a third: ``make coffee'' plus ``clear the
table'' composed into a morning routine. Composition requires shared interfaces and a
calculus over pre-/post-conditions; today's motion packages are opaque and do not compose,
and code skills compose only within a single authoring system.

\paragraph{Portability standards and skill ontologies.}
Reuse at scale needs a shared vocabulary: declared pre-conditions, expected effects, and
embodiment-capability profiles, with relations such as \emph{requires}, \emph{provides}, and
\emph{substitutable-by}. Prior work on skill ontologies and hardware-level reusability
points the way, but no cross-vendor standard exists, which is what a genuine ecosystem would
require.

\paragraph{Local versus shared libraries.}
Finally, the field must reconcile two notions of ``skill library'' that this survey has kept
distinct: the \emph{self-built} library a robot grows from its own experience (ASPIRE,
\S\ref{sec:code}) and the \emph{shared} library a robot downloads from a marketplace
(UniStore). The interesting regime is the hybrid---an agent that both contributes to and
draws from a commons---and its incentives, quality control, and feedback dynamics are
entirely unstudied.
